% callout-title-gothic.tex ----------------------------------------------------
% Quarto の callout は LaTeX では tcolorbox になり、タイトルは
%
%   title=\textcolor{...}{\faInfo}\hspace{0.5em}{注意すべき点}
%
% のように「書体指定なし」で埋め込まれます。全 callout のタイトル書体を
% まとめて変えるフックが tcolorbox の fonttitle です。
%
%   format:
%     pdf:
%       include-in-header: callout-title-gothic.tex
%
% ■ \AtBeginDocument が必要な理由
%   Quarto は include-in-header の内容を吐いたあとに tcolorbox を読み込むため、
%   プリアンブルに \tcbset を直書きすると
%     ! Undefined control sequence.  l.111 \tcbset
%   になります。\AtBeginDocument に包めば読み込み後に実行されます。
%
% ■ \bfseries だけでは足りない理由
%   \bfseries は「今の和文ファミリの太字」を選ぶので、本文が明朝なら
%   太明朝（NotoSerifCJKjp-Bold など）になります。ゴシックにするには
%   ファミリを明示的に切り替える必要があります。
%
%   \gtfamily … luatexja / jlreq / ltjsclasses / bxjs 系で使えるゴシック切替
%   \sffamily … fontspec や xeCJK では sansfont / CJKsansfont に切り替わる
%
%   下の \providecommand は、\gtfamily が無い構成のときだけ \sffamily を
%   代用するためのものです。どちらの構成でもそのまま動きます。
%   (fontspec 系では sansfont: に和文ゴシックを指定しておいてください。)
%------------------------------------------------------------------------------

\AtBeginDocument{%
  \providecommand{\gtfamily}{\sffamily}%
  \tcbset{fonttitle=\gtfamily\bfseries}%
}
